\section{Pengenalan Lingkungan Development vs Produksi}

Dalam pengembangan aplikasi web, terdapat beberapa \textbf{lingkungan} (\textit{environment}) yang masing-masing memiliki tujuan dan konfigurasi berbeda. \textbf{Development} (lokal) adalah tempat developer menulis dan menguji kode --- error ditampilkan, debugging aktif, data bersifat dummy. \textbf{Staging} adalah replika produksi untuk pengujian akhir sebelum rilis. \textbf{Production} (produksi) adalah lingkungan yang diakses pengguna nyata --- prioritas utama adalah keamanan, stabilitas, dan performa. Kesalahan konfigurasi di produksi dapat mengakibatkan kebocoran data atau downtime \cite{mdnDeploy}.

\begin{table}[ht]
\centering
\begin{tabular}{|P{3cm}|P{4cm}|P{4cm}|}
\hline
\textbf{Aspek} & \textbf{Development} & \textbf{Production} \\
\hline
\code{display\_errors} & On (tampilkan error) & Off (sembunyikan) \\
\hline
\code{error\allowbreak\_reporting} & E\_ALL & E\_ALL (log saja) \\
\hline
Database & Lokal, data dummy & Server terpisah, data asli \\
\hline
HTTPS & Tidak wajib & Wajib \\
\hline
Debug tools & Aktif & Nonaktif \\
\hline
Password & Sederhana (root/root) & Kuat dan unik \\
\hline
Backup & Tidak perlu & Rutin dan otomatis \\
\hline
\end{tabular}
\caption{Perbedaan lingkungan development dan production}
\end{table}

\begin{catatan}
Jangan pernah menggunakan konfigurasi development di produksi. Kesalahan umum: lupa mematikan \code{display\_errors}, menggunakan password database default, atau meninggalkan file debug (phpinfo.php) di server. Gunakan environment variable untuk membedakan konfigurasi per lingkungan.
\end{catatan}

